\documentclass[12pt,a4paper]{article}
\usepackage[a4paper,margin=20mm]{geometry}
\usepackage[T2A]{fontenc}
\usepackage[utf8]{inputenc}
\usepackage[russian]{babel}
\usepackage{amsmath,amssymb,mathtools}
\usepackage{array,booktabs,longtable}
\usepackage{xcolor}
\usepackage{hyperref}
\usepackage{tikz}

\hypersetup{
  colorlinks=true,
  linkcolor=blue!55!black,
  urlcolor=blue!55!black,
  pdftitle={Ревью домашней работы},
  pdfauthor={Лёвин Артём Александрович}
}

\setlength{\parindent}{0pt}
\setlength{\parskip}{0.55em}
\setlength{\tabcolsep}{3pt}
\renewcommand{\arraystretch}{1.22}

\begin{document}
\hypersetup{pageanchor=false}

\begin{titlepage}
\thispagestyle{empty}
\begin{center}
\vspace*{0.22\textheight}
{\Huge\bfseries Ревью домашней работы\par}
\vspace{2cm}
{\large Автор: Лёвин Артём Александрович\par}
\vspace{0.8cm}
{\large 14 сентября 2026 года\par}
\vfill
\begin{tikzpicture}
  \draw[blue!35!black,line width=0.8pt]
    (0,0) .. controls (2.2,0.45) and (4.4,-0.45) .. (6.6,0);
\end{tikzpicture}
\end{center}
\end{titlepage}
\hypersetup{pageanchor=true}

\newpage
\section*{Сводная таблица}
\small
\begin{tabular}{@{}>{\centering\arraybackslash}p{0.07\linewidth}
                  p{0.18\linewidth}
                  p{0.25\linewidth}
                  p{0.16\linewidth}
                  p{0.24\linewidth}@{}}
\toprule
\textbf{№ задания} & \textbf{Тема} & \textbf{Суть ошибки} & \textbf{Ваш ответ} & \textbf{Правильный ответ} \\
\midrule
\hyperref[task:2]{2} &
Действия с алгебраическими дробями &
При переписывании последней дроби числитель $28-a^2$ был заменён на $28a-a^2$. &
$5$ &
$\displaystyle \frac{3(6-a)}{5a-2}$, при $a\ne -3,\;\frac25,\;5,\;\frac{2}{11}$ \\
\addlinespace
\hyperref[task:3]{3} &
Общий знаменатель алгебраических дробей &
Ограничения указаны верно, но преобразование выражения не доведено до результата. &
не указан &
$\displaystyle \frac{3(x-1)}{x}$, при $x\ne0,\;-1$ \\
\bottomrule
\end{tabular}
\normalsize

\newpage
\thispagestyle{empty}
\vspace*{\fill}
\begin{center}
{\Huge\bfseries Разбор ошибок}
\end{center}
\vspace*{\fill}

\newpage
\phantomsection\label{task:2}
\section*{Задание 2}

\subsection*{Текст задания}
Упростить выражение
\[
\left(
\frac{36a^2}{5a^2+13a-6}
-
\frac{5a-2}{a+3}
\right)
:
\frac{11a-2}{a^2-2a-15}
-
\frac{28-a^2}{2-5a}.
\]

\subsection*{Ваше решение}
В первой части решения знаменатель
\[
5a^2+13a-6
\]
был разложен на множители, после чего первая разность была приведена к общему знаменателю. Получено
\[
\frac{36a^2-(5a-2)^2}{(5a-2)(a+3)}
=
\frac{11a^2+20a-4}{(5a-2)(a+3)}.
\]
Далее числитель был разложен как
\[
11a^2+20a-4=(11a-2)(a+2),
\]
а знаменатель второй дроби как
\[
a^2-2a-15=(a-5)(a+3).
\]
После деления на дробь и сокращения получен корректный промежуточный результат
\[
\frac{(a+2)(a-5)}{5a-2}.
\]
Затем в последней дроби вместо исходного числителя $28-a^2$ появилось выражение $28a-a^2$, после чего был получен ответ $5$.

\subsection*{Ваш ответ}
\[
5.
\]

\subsection*{Ошибка}
\textcolor{red}{Ошибка:} при переписывании последней дроби число $28$ было ошибочно заменено на произведение $28a$.

\subsection*{Где именно возникает ошибка}
Последний достоверно правильный шаг:
\[
\frac{(a+2)(a-5)}{5a-2}
-
\frac{28-a^2}{2-5a}.
\]
Первый неправильный шаг появляется при записи последней дроби в виде
\[
-\frac{28a-a^2}{2-5a}.
\]
В исходном выражении множителя $a$ после числа $28$ нет.

\subsection*{Почему этот шаг неверен}
При алгебраических преобразованиях можно изменять вид выражения только с помощью равносильных операций: раскрытия скобок, вынесения общего множителя, разложения на множители, приведения к общему знаменателю и других допустимых преобразований. Простое добавление множителя $a$ к одному из слагаемых не является равносильной операцией.

Выражения $28-a^2$ и $28a-a^2$ в общем случае различны. Например, при допустимом значении $a=0$ первое выражение равно $28$, а второе равно $0$. Поэтому после такой замены решается уже другое выражение, и дальнейшие правильные действия не могут восстановить исходную задачу.

\textcolor{blue!60!black}{Ключевая идея:} после каждого большого преобразования полезно сверять переписанный фрагмент с предыдущей строкой, особенно числители и знаки перед дробями. Здесь вся первая часть решения была выполнена правильно, а ошибка возникла именно при переносе последнего числителя.

\subsection*{Исправление от последнего правильного шага}
Начинаем с последней правильной строки:
\[
\frac{(a+2)(a-5)}{5a-2}
-
\frac{28-a^2}{2-5a}.
\]
Так как
\[
2-5a=-(5a-2),
\]
то
\[
\frac{28-a^2}{2-5a}
=
\frac{a^2-28}{5a-2}.
\]
Следовательно,
\[
\frac{(a+2)(a-5)}{5a-2}
-
\frac{a^2-28}{5a-2}
=
\frac{(a+2)(a-5)-(a^2-28)}{5a-2}.
\]
Раскрываем скобки в числителе:
\[
(a+2)(a-5)=a^2-3a-10.
\]
Поэтому
\[
\frac{a^2-3a-10-a^2+28}{5a-2}
=
\frac{18-3a}{5a-2}
=
\frac{3(6-a)}{5a-2}.
\]

\subsection*{Полное правильное решение}
Сначала фиксируем ограничения исходного выражения. Знаменатели не должны обращаться в нуль, а дробь, на которую выполняется деление, сама не должна быть равна нулю.

Разложим знаменатели:
\[
5a^2+13a-6=(5a-2)(a+3),
\]
\[
a^2-2a-15=(a-5)(a+3),
\]
\[
2-5a=-(5a-2).
\]
Отсюда
\[
a\ne -3,\qquad a\ne \frac25,\qquad a\ne5.
\]
Кроме того, делитель
\[
\frac{11a-2}{a^2-2a-15}
\]
не должен быть равен нулю, поэтому
\[
11a-2\ne0,
\qquad
 a\ne\frac{2}{11}.
\]
Итак, для исходного выражения
\[
a\ne -3,\quad \frac25,\quad 5,\quad \frac{2}{11}.
\]

Теперь упростим выражение. В скобках:
\begin{align*}
\frac{36a^2}{(5a-2)(a+3)}
-
\frac{5a-2}{a+3}
&=
\frac{36a^2-(5a-2)^2}{(5a-2)(a+3)}\\
&=
\frac{36a^2-(25a^2-20a+4)}{(5a-2)(a+3)}\\
&=
\frac{11a^2+20a-4}{(5a-2)(a+3)}\\
&=
\frac{(11a-2)(a+2)}{(5a-2)(a+3)}.
\end{align*}

Деление на дробь заменяем умножением на обратную дробь:
\begin{align*}
\frac{(11a-2)(a+2)}{(5a-2)(a+3)}
:
\frac{11a-2}{(a-5)(a+3)}
&=
\frac{(11a-2)(a+2)}{(5a-2)(a+3)}
\cdot
\frac{(a-5)(a+3)}{11a-2}\\
&=
\frac{(a+2)(a-5)}{5a-2}.
\end{align*}

Теперь обрабатываем последнюю дробь без изменения её числителя:
\[
\frac{28-a^2}{2-5a}
=
\frac{a^2-28}{5a-2}.
\]
Поэтому
\begin{align*}
\frac{(a+2)(a-5)}{5a-2}
-
\frac{a^2-28}{5a-2}
&=
\frac{a^2-3a-10-a^2+28}{5a-2}\\
&=
\frac{18-3a}{5a-2}\\
&=
\frac{3(6-a)}{5a-2}.
\end{align*}

\subsection*{Проверка}
Возьмём допустимое значение $a=0$.

Исходное выражение:
\[
\left(0-\frac{-2}{3}\right):\frac{-2}{-15}-\frac{28}{2}
=
\frac23:\frac{2}{15}-14
=
5-14
=-9.
\]
Полученная формула даёт
\[
\frac{3(6-0)}{5\cdot0-2}
=
\frac{18}{-2}
=-9.
\]
Значения совпали.

\textcolor{teal!60!black}{Итог:}
\[
\boxed{\displaystyle \frac{3(6-a)}{5a-2}},
\qquad
 a\ne -3,\;\frac25,\;5,\;\frac{2}{11}.
\]

\subsection*{Задача на закрепление}
Упростите выражение и укажите допустимые значения переменной:
\[
\left(
\frac{16b^2}{3b^2+10b-8}
-
\frac{3b-2}{b+4}
\right)
:
\frac{7b-2}{b^2-b-20}
-
\frac{18-b^2}{2-3b}.
\]

\newpage
\phantomsection\label{task:3}
\section*{Задание 3}

\subsection*{Текст задания}
Упростить выражение
\[
\frac{3x}{x+1}
+
\frac{2}{x}
-
\frac{2x+5}{x^2+x}.
\]

\subsection*{Ваше решение}
Вы верно выписали ограничения
\[
x\ne0,
\qquad
x\ne-1,
\]
и переписали исходное выражение. Дальнейшее приведение дробей к общему знаменателю и итоговое упрощение в работе отсутствуют.

\subsection*{Ваш ответ}
Не указан.

\subsection*{Ошибка}
\textcolor{red}{Ошибка:} решение не доведено до результата.

\subsection*{Где именно возникает ошибка}
Последний правильный шаг --- указание ограничений
\[
x\ne0,
\qquad
x\ne-1.
\]
Первого неверного преобразования здесь нет. Проблема в том, что после правильного начала решение остановлено: не выбран общий знаменатель, не объединены дроби и не получен итоговый результат.

\subsection*{Почему это неверно}
Ограничения показывают, при каких значениях переменной исходное выражение имеет смысл, но сами по себе не являются упрощением выражения. Чтобы выполнить задание полностью, нужно привести все три дроби к общему знаменателю, выполнить действия в числителе, разложить полученный числитель на множители и только затем сократить общий множитель, если он появляется.

\textcolor{blue!60!black}{Ключевая идея:} знаменатель третьей дроби сразу раскладывается как
\[
x^2+x=x(x+1).
\]
Это уже общий знаменатель для всех трёх дробей.

\subsection*{Исправление от последнего правильного шага}
Ограничения уже указаны верно:
\[
x\ne0,
\qquad
x\ne-1.
\]
Используем общий знаменатель $x(x+1)$:
\[
\frac{3x}{x+1}
=
\frac{3x^2}{x(x+1)},
\qquad
\frac{2}{x}
=
\frac{2(x+1)}{x(x+1)}.
\]
Третья дробь уже имеет нужный знаменатель, так как
\[
x^2+x=x(x+1).
\]
Получаем
\begin{align*}
\frac{3x^2}{x(x+1)}
+
\frac{2(x+1)}{x(x+1)}
-
\frac{2x+5}{x(x+1)}
&=
\frac{3x^2+2x+2-2x-5}{x(x+1)}\\
&=
\frac{3x^2-3}{x(x+1)}\\
&=
\frac{3(x-1)(x+1)}{x(x+1)}\\
&=
\frac{3(x-1)}{x}.
\end{align*}
При сокращении множителя $x+1$ исходное ограничение $x\ne-1$ сохраняется.

\subsection*{Полное правильное решение}
Начинаем с области допустимых значений. Из знаменателей
\[
x+1,\qquad x,\qquad x^2+x=x(x+1)
\]
получаем
\[
x\ne0,
\qquad
x\ne-1.
\]

Общий знаменатель равен
\[
x(x+1).
\]
Приводим каждую дробь к нему:
\[
\frac{3x}{x+1}
=
\frac{3x^2}{x(x+1)},
\]
\[
\frac{2}{x}
=
\frac{2(x+1)}{x(x+1)},
\]
\[
\frac{2x+5}{x^2+x}
=
\frac{2x+5}{x(x+1)}.
\]
Объединяем числители с учётом знака минус перед третьей дробью:
\begin{align*}
\frac{3x^2+2(x+1)-(2x+5)}{x(x+1)}
&=
\frac{3x^2+2x+2-2x-5}{x(x+1)}\\
&=
\frac{3x^2-3}{x(x+1)}\\
&=
\frac{3(x^2-1)}{x(x+1)}\\
&=
\frac{3(x-1)(x+1)}{x(x+1)}\\
&=
\frac{3(x-1)}{x}.
\end{align*}
Ограничения исходного выражения остаются
\[
x\ne0,
\qquad
x\ne-1.
\]

\subsection*{Проверка}
Возьмём допустимое значение $x=1$.

Исходное выражение:
\[
\frac{3}{2}+2-\frac{7}{2}=0.
\]
Упрощённое выражение:
\[
\frac{3(1-1)}{1}=0.
\]
Результаты совпадают.

\textcolor{teal!60!black}{Итог:}
\[
\boxed{\displaystyle \frac{3(x-1)}{x}},
\qquad
x\ne0,\;-1.
\]

\subsection*{Задача на закрепление}
Упростите выражение и укажите допустимые значения переменной:
\[
\frac{4y}{y-2}
+
\frac{3}{y}
-
\frac{7y+2}{y^2-2y}.
\]

\vfill
\begin{center}
\small Лёвин Артём Александрович эксклюзивно для Кристины
\end{center}

% Краткое сообщение Кристине: в задании 2 основная ошибка возникла при переписывании последней дроби — число двадцать восемь превратилось в произведение двадцати восьми на переменную. В задании 3 ограничения указаны верно, но решение не доведено до общего знаменателя и итогового упрощения. Обратите внимание прежде всего на точность переписывания выражений и на завершение преобразований до ответа.

\end{document}
